\documentclass[11pt]{article}
\usepackage[margin=1in]{geometry}
\usepackage{amsmath,amssymb,booktabs,longtable,hyperref}
\usepackage{listings}
\usepackage{xcolor}
\lstset{basicstyle=\ttfamily\small,breaklines=true,frame=single,columns=fullflexible}
\title{Independent Replication of \emph{Quantum Counting}\\
       (Brassard, H\o{}yer, Tapp, 1998)\\
       arXiv:quant-ph/9805082}
\author{Kukla \\ X-100 Replication Project, wave 2026-07-06}
\date{2026-07-06}

\begin{document}
\maketitle

\begin{abstract}
We independently reproduce the principal quantitative claim of Brassard,
H\o{}yer and Tapp's 1998 paper \emph{Quantum Counting}: Theorem~5's error bound
$|t-\tilde t| < \tfrac{2\pi}{P}\sqrt{tN} + \tfrac{\pi^2}{P^2}N$ with success
probability $\ge 8/\pi^2$ for the counting algorithm \textsc{Count}$(F,P)$ that
combines a Grover iteration with quantum phase estimation. Two independent
Python simulators were built --- an exact analytic QPE marginal (via the
Dirichlet kernel on the two Grover eigenphases) and a full gate-level Qiskit~2.5
statevector circuit --- and cross-validated to $\|\cdot\|_\infty \le 3\!\times\!10^{-15}$
across seven diverse $(n,t,p)$ cases. A 90-configuration sweep over
$N\in\{16,32,64\}$, varying~$t$ and counting qubits $p\in\{3,\dots,8\}$, satisfies
Theorem~5's bound and success-probability guarantee in $90/90$ cases. LLM-judge
(Argo \texttt{argo:gpt-5.4}) verdict: \textbf{PARTIAL} (agreement 1.0,
coverage 0.67): the counting theorem and its counting-qubit scaling are
strongly reproduced; the broader amplitude-amplification claims (Theorems 1--3)
were not exercised numerically.
\end{abstract}

\section{Paper summary}
The 1998 paper unifies Grover's algorithm and Shor's phase estimation into a
\emph{quantum counting} algorithm. Given a Boolean oracle
$F:\{0,\ldots,N-1\}\to\{0,1\}$ with $t$ marked inputs and a search space
$N=2^n$, the algorithm \textsc{Count}$(F,P)$ uses $P=2^p$ counting qubits to
apply controlled powers of the Grover operator
\begin{equation}
G = (2|\psi\rangle\langle\psi| - I)\,O_F, \qquad
|\psi\rangle = H^{\otimes n}|0\rangle,
\end{equation}
and inverse quantum Fourier transform on the counting register. The two
non-trivial eigenphases of $G$ are $\pm 2\theta$ with
$\sin^2\theta = t/N$; measuring the counting register yields an integer $f$ from
which $\tilde t = N\sin^2(\pi f / P)$ is computed. The paper proves
\begin{equation}
|t - \tilde t| < \frac{2\pi}{P}\sqrt{tN} + \frac{\pi^2}{P^2}\,N
\qquad\text{with probability}\ \ge \frac{8}{\pi^2}\approx 0.811.
\label{eq:thm5}
\end{equation}

\section{Claims table}
\begin{center}
\begin{tabular}{@{}llll@{}}
\toprule
\textbf{ID} & \textbf{Claim} & \textbf{Testable?} & \textbf{Tested?}\\
\midrule
C1 & Amplitude amplification: $\Theta(1/\sqrt{a})$ speedup for any $A$ (Thm~1--3) & Yes (asymptotic) & No\\
C2 & \textsc{Count}$(F,P)$ satisfies Eq.~(\ref{eq:thm5}) with prob.\ $\ge 8/\pi^2$ (Thm~5/6) & Yes (numerical) & \textbf{Yes}\\
C3 & Corollary~2: $P=c\sqrt N \Rightarrow |t-\tilde t| < \tfrac{2\pi}{c}\sqrt t + \tfrac{\pi^2}{c^2}$ & Yes (numerical) & Yes (consistency)\\
C4 & Heuristic-search quadratic speedup (Thm~4) & Structural / setup-specific & No\\
\bottomrule
\end{tabular}
\end{center}

\section{Method}
\begin{enumerate}
  \item \textbf{Paper retrieval.} PDF pulled from
        \url{https://arxiv.org/pdf/quant-ph/9805082}
        (176\,410 bytes, 12 pages, PDF~1.4).
  \item \textbf{Extraction.} \texttt{pdftotext -layout} into
        \texttt{work/paper.txt} (581 lines), copied into
        \texttt{extraction/marker.md} and \texttt{extraction/nougat.mmd} with
        header noting the pdftotext substitution (Marker/Nougat not installed
        locally and paper not in central corpus).
  \item \textbf{Environment.} Python venv at \texttt{work/venv},
        Qiskit~2.5.0 + qiskit-aer~0.17.2 + NumPy.
  \item \textbf{Analytic implementation
        (\texttt{work/quantum\_counting.py}).}
        For each $(n, \text{marked}, p)$ we compute the exact marginal
        distribution over the measured integer $f$ using the standard QPE
        Dirichlet-kernel form:
        \[
          P(f\mid\phi) = \frac{\sin^2(\pi P(\phi - f/P))}{P^2\sin^2(\pi(\phi - f/P))}.
        \]
        The state $|\psi\rangle$ decomposes as an equal superposition of the
        two eigenvectors of $G$ with eigenphases
        $\phi_+ = \theta/\pi$ and $\phi_- = 1 - \theta/\pi$, so the counting
        register's marginal is $\tfrac12 P(f\mid\phi_+) + \tfrac12 P(f\mid\phi_-)$.
  \item \textbf{Gate-level cross-check (\texttt{work/verify\_qiskit.py},
        \texttt{work/verify\_qiskit\_multi.py}).}
        A Qiskit circuit with $p$ counting qubits and $n$ search qubits, initial
        Hadamards, controlled-$G^{2^j}$ (control $=$ \texttt{count\_reg}$[j]$),
        and inverse QFT with swaps is built. The full statevector is computed
        and marginalised to the counting register. Convention debugging (see
        \S\ref{sec:worked-didnt}) uncovered a subtle Qiskit little-endian issue
        with the control qubit position in \texttt{UnitaryGate.append}: to use
        the block-diagonal $\mathrm{diag}(I,U)$ form, the control must be the
        \emph{last} argument (MSB of the gate's matrix index). After the fix,
        analytic and gate-level marginals agree to
        $\|\cdot\|_\infty \le 3\!\times\!10^{-15}$ across seven diverse cases.
  \item \textbf{Sweep (\texttt{work/quantum\_counting.py}).}
        Configurations: $n\in\{4,5,6\}$, $t$ values per~$n$ ($\{1,3,6,10,15\}$;
        $\{1,5,12,20,27,31\}$; $\{1,10,25,40,55,63\}$), and
        $p\in\{3,4,5,6,7,8\}$, giving 90 rows. For each row we record
        $\tilde t_{\mathrm{argmax}}$, $|t-\tilde t|$, the Theorem~5 bound, and
        the \emph{exact} probability that the measurement lands in the
        ``success'' set $\{f : |t - N\sin^2(\pi f/P)| < \text{bound}\}$.
  \item \textbf{LLM-judge (\texttt{work/llm\_judge.py}).}
        Free Argo endpoint (\texttt{localhost:44497}, \texttt{argo:gpt-5.4}).
        Prompt included the sweep summary, cross-verification log, and paper
        text; judge returns per-claim verdicts + overall JSON.
        Raw output at \texttt{report/evidence/llm\_judge\_raw.txt}; parsed at
        \texttt{report/evidence/llm\_judge.json}.
\end{enumerate}

\section{Results vs paper}
\subsection{Sanity: gate-level vs analytic (7 cases)}
Worst $\|\cdot\|_\infty$ deviation: $2.72\times10^{-15}$. All argmax pairs are
either identical or mirror-symmetric ($f$ vs $P-f$), which is expected because
QPE assigns equal probability to $\phi_+ = \theta/\pi$ and
$\phi_- = 1 - \theta/\pi$, and both map to the same
$\tilde t = N\sin^2(\pi f/P)$.
See \texttt{report/evidence/qiskit\_verify\_multi.log}.

\subsection{Main sweep (90 configs)}
Every one of the 90 configurations satisfied both the Theorem~5 error bound and
the $8/\pi^2$ success-probability guarantee:
\begin{center}\begin{tabular}{lc}\toprule
Argmax within Theorem-5 bound & 90 / 90\\
Exact $P(\text{within bound}) \ge 8/\pi^2$ & 90 / 90\\
\bottomrule\end{tabular}\end{center}

\noindent Representative rows ($t$ = true count, $\tilde t$ = argmax estimate,
$B$ = Thm-5 bound, $p_s$ = exact success prob):
\begin{center}\small\begin{tabular}{rrrrrrr}\toprule
$n$ & $t$ & $p$ & $\tilde t$ & $|t-\tilde t|$ & $B$ & $p_s$\\
\midrule
4 & 1 & 8 & 1.039 & 0.039 & 0.101 & 1.000\\
4 & 3 & 6 & 2.925 & 0.075 & 0.719 & 0.996\\
5 & 12 & 6 & 11.36 & 0.645 & 2.00 & 1.000\\
5 & 20 & 7 & 19.89 & 0.112 & 1.26 & 1.000\\
6 & 25 & 8 & 24.99 & 0.011 & 0.99 & 1.000\\
6 & 40 & 5 & 38.24 & 1.76 & 10.55 & 1.000\\
6 & 63 & 4 & 61.56 & 1.44 & 27.40 & 1.000\\
\bottomrule\end{tabular}\end{center}
The error monotonically shrinks with~$P$ as $\sim P^{-1}$ in the leading term,
consistent with Corollary~2.

\section{Verdict}
\textbf{PARTIAL} (with agreement $=1.00$ on the tested claims, coverage $=0.67$).
Justification: the paper's principal quantitative counting result (C2/Theorem~5)
is reproduced independently at machine precision, in every one of 90 test
configurations across three search-space sizes. Corollary~2's scaling (C3) is
consistent with our sweep. The general amplitude-amplification speedup claims
(C1) and heuristic-search theorem (C4) were not exercised numerically in this
replication and would require additional benchmarks with families of algorithms
$A$ of varying success probability~$a$.

\section{What worked and what didn't}\label{sec:worked-didnt}
\subsection*{Worked}
\begin{itemize}
  \item Analytic QPE marginal derivation (Dirichlet kernel over the two Grover
        eigenphases) --- gave a closed-form target the gate-level implementation
        could be checked against, avoiding a fully-blind reproduction.
  \item Exact success-probability computation --- we sum the marginal over the
        theorem's success set instead of relying on shot statistics, giving the
        crispest possible check of the $8/\pi^2$ bound.
  \item Cross-validation to $10^{-15}$ ruled out both circuit and analytic bugs.
\end{itemize}
\subsection*{Didn't work at first}
\begin{itemize}
  \item First Qiskit gate-level attempt disagreed with the analytic form by
        $\|\cdot\|_\infty \approx 0.6$ and returned nonsense $\tilde t$ values.
        Root cause: with the standard block form $\mathrm{full}=\mathrm{diag}(I,U)$
        passed to \texttt{UnitaryGate}, the control qubit must be the
        \emph{last} argument in \texttt{qc.append(cU, [\dots, control])} in
        Qiskit's little-endian convention (MSB of the gate matrix index).
        Placing the control as the first argument silently applied~$U$ to a
        wrong superposition of control states. Fix: swap to
        \texttt{list(search\_reg) + [count\_reg[j]]}.
  \item Argo \texttt{claude-opus-4.8} judge call failed with an upstream JSON
        parse error on the large prompt. Fell back to \texttt{argo:gpt-5.4},
        which returned a well-formed judgement.
\end{itemize}
\subsection*{Caveats}
\begin{itemize}
  \item The sweep is restricted to $n\le 6$ (dense-matrix $G$ powers become
        expensive beyond that). The theorem is asymptotic in nature; our
        agreement in the small-$N$ regime is a necessary but not fully
        exhaustive test.
  \item Corollary~2's exact scaling $P=c\sqrt N$ was not \emph{fit}
        parametrically; we only verified that shrinking $P$-dependent bound
        held pointwise.
\end{itemize}

\section{Open Questions}
\begin{description}
  \item[Q1.] \textbf{Tightness of the $8/\pi^2$ probability floor for
    intermediate $t/N$.} The exact success probability in our sweep was
    typically $\ge 0.99$ across all $(n,t,p)$, well above the theoretical
    $\approx 0.811$. Where does the paper's bound bite? For which $t/N$ and $P$
    is the marginal actually near the floor?
  \item[Q2.] \textbf{Behaviour when $t/N \to 0$ or $t/N \to 1$.} The eigenphase
    $\theta \to 0$ regime gives progressively worse leading-term error
    $\propto\sqrt{tN}/P$; empirically we saw the argmax hop to nearby $f$ but
    remained within the bound. What is the actual worst-case
    $t/N$ for a fixed $P$?
  \item[Q3.] \textbf{Bias of $\tilde t_{\mathrm{argmax}}$ vs the paper's
    ``round'' rule.} The paper does not commit to argmax; it just bounds any
    output of the algorithm. In practice \emph{which} decoding rule (argmax,
    weighted median, matched-filter over the two peaks) is best for a fixed
    query budget?
  \item[Q4.] \textbf{Non-uniform initial states.} The whole analysis assumes
    the search register starts in the exact uniform superposition $A|0\rangle =
    H^{\otimes n}|0\rangle$. What happens under imperfect $A$ (e.g.\
    depolarising or dephasing noise on the counting or search register), which
    is inevitable on real hardware?
  \item[Q5.] \textbf{Comparison to iterative-QPE / Bayesian counting.} Modern
    variants (Grinko--et-al.\ iterative amplitude estimation; Bayesian phase
    estimation) claim comparable or better sample complexity with fewer
    counting qubits. Under our exact framework, at what $N,t$ does one method
    dominate?
\end{description}

\end{document}
